\begin{proposition}[5.6.4]
\label{IV.2.s5.6.4-fr}
Todo anillo cociente de un anillo regular es universalmente catenario.
\end{proposition}

Todo se reduce a ver que un anillo regular $A$ es universalmente catenario
(5.6.3, (iv)); ahora bien, se sabe que $A[T_1,\ldots,T_n]$ es regular (\textbf 0, 17.3.7),
por tanto catenario (\textbf 0, 16.5.12), y se concluye por (5.6.1, \emph{a)}).

En particular, se sigue del teorema de Cohen (\textbf 0, 19.8.8) que todo \emph{anillo
local noetheriano completo} es universalmente catenario. Del mismo modo, como
toda álgebra de tipo finito sobre un cuerpo es cociente de un anillo regular,
todo preesquema \emph{localmente de tipo finito sobre un cuerpo} es universalmente
catenario.

Veremos más adelante (6.3.7) que en (5.6.4) se puede sustituir « anillo
regular » por « anillo de Cohen-Macaulay ».

\begin{proposition}[5.6.5]
\EGAReaderTarget{ega4.c04.prop.5-6-5}{IV.2.s5.6.5-fr}
Sean $Y$ un preesquema irreducible localmente noetheriano, $X$ un preesquema
irreducible, $f:X\to Y$ un morfismo dominante localmente de tipo finito. Sea
$\xi$ \textup{(resp. $\eta$)} el punto genérico de $X$ \textup{(resp. $Y$)}, y sea
$e=\dim(f^{-1}(\eta))=\operatorname{deg.tr}_{\mathbf{k}(\eta)}\boldsymbol{k}(\xi))$
\textup{(« dimensión de la fibra genérica », cf. $(\mathbf{0}_{\mathrm I},2.1.8)$ y
(4.1.1))}. Para todo $x\in X$, se tiene entonces, poniendo $y=f(x)$,
\begin{equation}
e+\dim(\mathcal O_y)
\geq\operatorname{deg.tr}_{\mathbf{k}(y)}\boldsymbol{k}(x)+\dim(\mathcal O_x)
\tag{5.6.5.1}\label{IV.2.s5.6.5.1-fr}
\end{equation}
\begin{equation}
\dim(\mathcal O_x)
\leq\dim(\mathcal O_y)
+\dim(\mathcal O_x\otimes_{\mathcal O_y}\boldsymbol{k}(y))\setmin\delta(x)
\tag{5.6.5.2}\label{IV.2.s5.6.5.2-fr}
\end{equation}
poniendo $\delta(x)=\dim_x(f^{-1}(y))\setmin e$.

Además, si $Y$ es universalmente catenario, los dos miembros de
\textup{(5.6.5.1) (resp. (5.6.5.2))} son iguales. Si, además, $x$ es cerrado en
$f^{-1}(y)$, se tiene
\begin{equation}
\dim(\mathcal O_x)=\dim(\mathcal O_y)+e.
\tag{5.6.5.3}\label{IV.2.s5.6.5.3-fr}
\end{equation}
\end{proposition}

Es posible evidentemente limitarse al caso en que $X$ e $Y$ son afines (por
tanto $f$ de tipo finito) y reducidos (1.5.4), por tanto íntegros. Como $f$ es
dominante, $\mathcal O_y$ se identifica entonces con un subanillo de
$\mathcal O_x$, y los cuerpos de fracciones respectivos de $\mathcal O_y$ y
de $\mathcal O_x$ con $\boldsymbol{k}(\eta)$ y $\boldsymbol{k}(\xi)$ (\textbf I, 8.2.7); además,
$\mathcal O_x$ es el anillo local de una $\mathcal O_y$-álgebra íntegra de
tipo finito $B$ en un ideal primo $\mathfrak q$ (\textbf I, 3.6.5);
aplicando
(5.5.8.1) a $A=\mathcal O_y$, $B$ y $\mathfrak q$, se obtiene (5.6.5.1).
Además,
$\mathcal O_x\otimes_{\mathcal O_y}\boldsymbol{k}(y)=\mathcal O_x/\mathfrak m_y\mathcal O_x$
no es sino el anillo local de la fibra $f^{-1}(y)$ en el punto $x$
(\textbf I, 3.6.1); como $f^{-1}(y)$ es un preesquema de tipo finito sobre $\boldsymbol{k}(y)$, se
sigue de (5.2.3) que se tiene
\[
\dim(\mathcal O_x\otimes_{\mathcal O_y}\boldsymbol{k}(y))
=\dim_x(f^{-1}(y))\setmin\operatorname{deg.tr}_{\mathbf{k}(y)}\boldsymbol{k}(x);
\]
la desigualdad (5.6.5.2) no es, pues, sino otra forma de (5.6.5.1).

El hecho de que los dos miembros de (5.6.5.1) sean iguales cuando $Y$ es
universalmente
\oldpage[IV-2]{100}
catenario no es sino la igualdad (5.6.1.1), aplicada a
$A=\mathcal O_y$ y $B'=\mathcal O_x$. Para demostrar que se tiene además
(5.6.5.3) cuando $x$ es cerrado en $f^{-1}(y)$, basta observar que, en
general, $\operatorname{deg.tr}_{\mathbf{k}(y)}\boldsymbol{k}(x)$ es la dimensión de
$\overline{\{x\}}\cap f^{-1}(y)$, según se sigue de (5.2.1) aplicado al
subpreesquema cerrado reducido de $f^{-1}(y)$ que tiene este subespacio como
espacio subyacente.

Demostraremos más adelante (13.1.1) que, bajo las condiciones de (5.6.5), \emph{se
tiene siempre} $\delta(x)\geq0$, y (5.6.5.2) precisa, pues, en este caso (5.5.2.1).
